\documentclass[a4paper, 12pt]{article}
\usepackage[a4paper, left=1.5cm, right=1.5cm, top=1.5cm, bottom=1.5cm]{geometry}

\usepackage{polyglossia}
\setmainlanguage{russian}
\setmainfont{Times New Roman}

\begin {document}
	\begin{center}
		\large План исследования по теме: "Конфликт интересов Российской и Британской империй в начале ХХ века"
	\end{center}
	
	\textbf{Аннотация: } Россию и Великобританию в начале XX века связывали сложные отношения: острое соперничество в центральной Азии, доставшееся в наследство от "Большой игры" конца XIX века; тесное сотрудничество: Англия после 1907 года являлась одним из основных инвесторов в Российскую промышленность, добычу полезных ископаемых; стратегическая поддержка: империи считали опасным рост силы Германии, поэтому они старались помешать строительству Багдадской железной дороги, которая создала бы проблемы для обеих. Поэтому я считаю эту тему весьма интересной для изучения, также на примере данных взаимоотношений можно будет лучше понять правила, по которым выстраиваются отношения между странами в современном мире.
	
	
	\textbf{\textit{План работ:}} 
	\begin{enumerate}
		\item Найти необходимый список источников, на котором будет основана работа: В электронных читальных залах можно изучить архивы: "Статистические сведения ввоза и вывоза товаров по странам за 1900-1914 годы", чтобы получить некоторые сведения о объеме торговли между странами, также необходимо будет проверить открытые источники в сети Интернет на темы уже озвученные в аннотации. \textbf{Срок выполнения: 15 марта}
		\item Составление сценария видеоролика, важно учесть, что длина видеоролика не более 10 минут. \textbf{Срок выполнения: 1 апреля}
		\item Подбор визуальных источников. Можно попробовать снять некоторые сцены в среде города, чтобы сделать ролик более живым. Необходимо выполнить раскадровку сценария. \textbf{Срок выполнения: 14 апреля} 
		\item Итоговый сбор и монтаж видеоролика. \textbf{Срок выполнения: 3 мая}
		
	\end{enumerate}
	
\end{document}